\section{Framework: intercept disagreement in an opinionated-markets model}
\label{sec:framework}

Caballero \& Simsek (2022) study a Fed and a market that disagree, each confident in its own
view, with the Fed setting policy on its belief and the market pricing assets knowing it will.
The market therefore anticipates policy ``mistakes'' and prices the resulting risk. Their state
variable is aggregate demand. This section runs the same logic with the disagreement located in
the \emph{intercept of the policy rule}, and the change of state variable turns out to matter:
demand disagreement mean-reverts, whereas disagreement about a constant does not, and that
difference generates predictions their model does not make.

\subsection{Setup}

Both parties know the rule and observe the state. They differ only in the neutral real rate:
the Fed believes $r^{*F}$, the market believes $r^{*m}$, and
\begin{equation}
\delta_t \;\equiv\; r^{*m}_t - r^{*F}_t .
\label{eq:delta}
\end{equation}
The Fed sets policy on its own estimate,
\begin{equation}
i_t \;=\; r^{*F}_t + \pi^{*} + \phi_\pi\!\left(\pi_t - \pi^{*}\right) + \phi_x x_t ,
\qquad \phi_\pi > 1 ,
\label{eq:rule}
\end{equation}
and the market evaluates the resulting stance against \emph{its} neutral rate, perceiving policy
to be $\delta_t$ looser than the Fed intends.

\subsection{The market's steady state is not the Fed's}

Suppose the disagreement persists and the economy settles. In any steady state the real rate
equals the market's neutral rate, $i - \pi = r^{*m}$. Substituting \eqref{eq:rule},
\[
r^{*m} \;=\; r^{*F} + \pi^{*} + \phi_\pi(\pi-\pi^{*}) - \pi
\qquad\Longrightarrow\qquad
\left(\phi_\pi-1\right)\left(\pi-\pi^{*}\right) \;=\; \delta ,
\]
so the market expects inflation to settle permanently away from target at
\begin{equation}
\boxed{\;\pi - \pi^{*} \;=\; \frac{\delta_t}{\phi_\pi - 1}\;}
\label{eq:bias}
\end{equation}
A Fed that underestimates the neutral rate holds policy too easy forever, and under a rule with
$\phi_\pi>1$ the only way the rule can be satisfied is with inflation permanently above target.
The macro half of this is Orphanides' account of the Great Inflation --- a mismeasured natural
rate producing a sustained inflation bias --- imported here as a belief the \emph{market} holds
about the Fed.

\subsection{The amplification result}

The nominal endpoint the market prices is $r^{*m} + \pi^{*} + \delta/(\phi_\pi-1)$; the endpoint
the FOMC publishes is $r^{*F} + \pi^{*}$. The difference is
\begin{equation}
\boxed{\;\underbrace{i^{*m}_t - i^{*F}_t}_{\text{priced endpoint gap}} \;=\; \frac{\phi_\pi}{\phi_\pi-1}\;\delta_t\;}
\label{eq:amplification}
\end{equation}
\emph{A small disagreement about $r^{*}$ implies a large disagreement about where nominal rates
settle}, with the multiplier set by the perceived aggressiveness of the rule:

\begin{center}\small
\begin{tabular}{lccc}
\toprule
$\phi_\pi$ & inflation bias per 1 pp of $\delta$ & endpoint gap per 1 pp of $\delta$ & endpoint gap when sd$(\delta)=0.105$\\ \midrule
1.25 & 4.00 & 5.00 & 52 bp\\
1.50 & 2.00 & 3.00 & 32 bp\\
2.00 & 1.00 & 2.00 & 21 bp\\
2.50 & 0.67 & 1.67 & 18 bp\\
\bottomrule\end{tabular}
\end{center}

\noindent This reverses the reading of Fact~2. The measured Fed-minus-dealer $r^{*}$ gap has a
standard deviation of only 10 bp, which looked fatal for any design built on $\delta$. Through
\eqref{eq:amplification} that same gap implies 18--52 bp of disagreement about the nominal
endpoint --- an order of magnitude that moves long yields. The smallness of the $r^{*}$ gap is a
prediction of the model, not an obstacle to it: disagreements about the intercept are arbitraged
down in the $r^{*}$ dimension precisely because they are amplified in the priced dimension.

\subsection{Pricing, with the maturity profile derived rather than imposed}

Let the disagreement decay as the Fed learns, $E_t[\delta_{t+j}] = (1-\mu)^{j}\delta_t$, and let
the inflation bias build with the Phillips-curve pass-through, so that the bias at horizon $j$ is
$\lambda_j \delta_t$ with
\begin{equation}
\lambda_j \;=\; \frac{1}{\phi_\pi-1}\underbrace{\left(1-e^{-\kappa j}\right)}_{\text{bias builds}}\underbrace{\left(1-\mu\right)^{j}}_{\text{Fed learns}} .
\label{eq:lambda}
\end{equation}
Averaging expected short rates then gives
\begin{equation}
y^{(n)}_t \;=\; \underbrace{r^{*m}_t + \pi^{*}}_{\text{market endpoint}}
\;+\; \underbrace{\omega(\rho,n)\,\text{gap}_t}_{\text{stance}}
\;+\; \underbrace{\Lambda(n)\,\delta_t}_{\text{perceived error}}
\;+\; \underbrace{TP^{(n)}\!\left(\sigma^{2}_{\delta,t}\right)}_{\text{disagreement premium}} ,
\qquad \Lambda(n)=\frac1n\sum_{j<n}\lambda_j .
\label{eq:pricing}
\end{equation}
$\Lambda(n)$ is now a derived object with a hump: it is near zero at short maturities because the
bias has not yet accumulated, peaks where the build-up and the learning decay cross, and falls
away as $(1-\mu)^j$ dominates. \textbf{The peak location is pinned by $\kappa$ and $\mu$ alone},
both estimable outside the event sample, so the hump is falsifiable in a way an asserted shape
is not.

\subsection{Predictions}

\begin{enumerate}\setlength\itemsep{0.25em}
\item \textbf{Amplification.} The gap between the market-implied nominal endpoint and the FOMC's
published dot is $\phi_\pi/(\phi_\pi-1)$ times the measured $r^{*}$ gap. Both sides are
observable, so the ratio \emph{identifies the market's perceived $\phi_\pi$} --- which can then
be compared against the survey-based estimates of Bauer, Pflueger \& Sunderam (2024). An
over-identifying restriction, not a free parameter.
\item \textbf{A composition restriction on the TIPS split.} The endpoint gap divides into
$\delta$ in the real dimension and $\delta/(\phi_\pi-1)$ in the inflation dimension, so
\[
\frac{\text{real endpoint gap}}{\text{breakeven gap}} \;=\; \phi_\pi - 1 .
\]
At $\phi_\pi=1.5$ breakevens should move twice as much as real forwards, \emph{in the same
direction}. This is sharper than a sign prediction and it is testable on observed nominal and
TIPS curves with no decomposition of either.
\item \textbf{Sign.} $\delta>0$ --- the Fed's $r^{*}$ below the market's --- means policy is
perceived too easy, so both the inflation bias and the endpoint gap are positive. $\delta<0$
reverses both.
\item \textbf{Maturity profile.} $\Lambda(n)$ is hump-shaped with a peak determined by $\kappa$
and $\mu$; the stance term $\omega(\rho,n)$ declines monotonically; the endpoint enters flat.
Three distinguishable shapes.
\item \textbf{State dependence, and a link to the perceived-slope literature.} The multiplier
$\phi_\pi/(\phi_\pi-1)$ is decreasing in $\phi_\pi$: \emph{the same $r^{*}$ disagreement matters
more when markets believe the Fed responds weakly}. Intercept disagreement and slope perception
are complements, not rival explanations, which is a testable interaction rather than a
positioning claim.
\item \textbf{The premium.} Uncertainty about $\delta$ propagates into the whole expected path
undamped by horizon, so term-premium sensitivity to $\sigma^2_\delta$ should rise with maturity
and not decay --- unlike stance uncertainty, which is damped by $\omega$.
\end{enumerate}

\subsection{What the model assumes, stated plainly}

\emph{Dogmatic disagreement.} Prediction 1 rests on the market expecting the Fed \emph{not} to
learn quickly. With fast learning ($\mu$ large) the steady state is never approached and the
amplification shrinks toward $\delta$ itself; $\mu$ is therefore the parameter the amplification
result lives or dies on, and it is estimable from the observed persistence of the Fed-dealer gap.

\emph{The Taylor principle.} Everything here requires $\phi_\pi>1$, and the multiplier diverges
as $\phi_\pi\to1$. That is not a defect --- it is the model saying intercept errors are
catastrophic for a weakly-responding central bank --- but it means the empirical work must
bound $\phi_\pi$ away from one, and estimates near the boundary should be treated as evidence
against the specification rather than as very large effects.

\emph{Whose endpoint anchors the yield.} \eqref{eq:pricing} anchors on $r^{*m}$, the market's
belief, which is the correct choice for a pricing equation and the opposite of the earlier draft.
The Fed's estimate enters only through the bias term. If instead one believes realised short
rates are generated by the Fed's rule regardless of the truth, the anchor is
$\theta r^{*m} + (1-\theta)r^{*F}$; $\theta$ is estimable from the maturity profile and nests
both readings.
